% Part: proof-theory
% Chapter: sequent-calculus
% Section: proof-examples

\documentclass[../../../include/open-logic-section]{subfiles}

\begin{document}

\olfileid{pt}{seq}{ex}

\olsection{Tuladha Bukti}

Kanggo menehi !!{proof}s sekuen, lumrahe luwih gampang makarya
munggah saka sekuen pungkasan: pilih !!a{formula} ing sekuen kasebut
minangka !!{formula} aktif, tulis premis aturan sing cocog ing ndhuwure,
banjur baleni nganti tekan aksioma. Upamane kita arep mbuktekake sekuen
\[(!C \land !D) \lif !E \Sequent !C \lif (!D \lif !E)\]
Gatekna $!C \lif (!D \lif !E)$: formula iki awujud $!A \lif !B$
lan ana ing sisih tengen sekuen. Nyocogake kabeh sekuen karo aturan
$\RightR{\lif}$ saka \Log{G1c},
\[\Axiom$ !A, \Gamma \fCenter \Delta, !B$
\RightLabel{\RightR{\lif}}
\UnaryInf$ \Gamma \fCenter \Delta, !A \lif !B $
\DisplayProof\]
menehi pilihan $\Gamma = \{(!C \land !D) \lif !E\}$,
$\Delta = \emptyset$, $!A \ident !C$, lan $!B \ident !D \lif !E$.
Dadi inferensi pungkasan bukti bisa awujud
\[\Axiom$ !C, (!C \land !D) \lif !E \fCenter !D \lif !E$
\RightLabel{\RightR{\lif}}
\UnaryInf$(!C \land !D) \lif !E \fCenter !C \lif (!D \lif !E)$
\DisplayProof\]
Yen gagasan iki dibaleni, kanthi $!D \lif !E$ minangka formula aktif,
kita oleh
\[\Axiom$ !D, !C, (!C \land !D) \lif !E \fCenter !E$
\RightLabel{\RightR{\lif}}
\UnaryInf$!C, (!C \land !D) \lif !E \fCenter !D \lif !E$
\RightLabel{\RightR{\lif}}
\UnaryInf$(!C \land !D) \lif !E \fCenter !C \lif (!D \lif !E)$
\DisplayProof\]
Saiki kita kudu nggatekake $(!C \land !D) \lif !E$ lan nggunakake
aturan $\LeftR{\lif}$ saka \Log{G1c},
\[\Axiom$ \Gamma \fCenter \Delta, !A$
\Axiom$ !B, \Gamma \fCenter \Delta$
\RightLabel{\LeftR{\lif}}
\BinaryInf$ !A \lif !B, \Gamma \fCenter \Delta$
\DisplayProof\]
Iki menehi
\[\Axiom$!D, !C \fCenter !E, !C \land !D$
\Axiom$!D, !C, !E \fCenter !E$
\RightLabel{\LeftR{\lif}}
\BinaryInf$ !D, !C, (!C \land !D) \lif !E \fCenter !E$
\RightLabel{\RightR{\lif}}
\UnaryInf$!C, (!C \land !D) \lif !E \fCenter !D \lif !E$
\RightLabel{\RightR{\lif}}
\UnaryInf$(!C \land !D) \lif !E \fCenter !C \lif (!D \lif !E)$
\DisplayProof\]
Kanthi !!{formula} $!C \land !D$ ing sekuen paling ndhuwur kiwa
minangka formula utama lan aturan $\RightR{\land}$, kita oleh
\[
\Axiom$!D, !C \fCenter !E, !C$
\Axiom$!D, !C \fCenter !E, !D$
\RightLabel{\RightR{\land}}
\BinaryInf$!D, !C \fCenter !E, !C \land !D$
\Axiom$!D, !C, !E \fCenter !E$
\RightLabel{\LeftR{\lif}}
\BinaryInf$ !D, !C, (!C \land !D) \lif !E \fCenter !E$
\RightLabel{\RightR{\lif}}
\UnaryInf$!C, (!C \land !D) \lif !E \fCenter !D \lif !E$
\RightLabel{\RightR{\lif}}
\UnaryInf$(!C \land !D) \lif !E \fCenter !C \lif (!D \lif !E)$
\DisplayProof
\]
Nganti kene kabeh bener. Gatekna yen kabeh langkah iki uga bisa
ditindakake ing \Log{G3c}, amarga aturan sing wis digunakake padha
ing \Log{G1c} lan~\Log{G3c}. Kita wis oleh pecahan !!a{proof}
sing saben sekuen paling ndhuwure ngemot !!{formula} sing padha
ing kiwa lan tengen. Nanging apa sekuen iki minangka aksioma
gumantung marang sisteme.

Ing \Log{G1c}, kita kudu menehi !!{proof}s kanggo sekuen paling
ndhuwur saka aksioma awujud $!A \Sequent !A$. Iki gampang ditindakake
nganggo aturan pelemahan. Upamane, sekuen paling ndhuwur kiwa
$!D, !C \Sequent !E, !C$ bisa !!{prove} kaya iki:
\[
\Axiom$!C \fCenter !C$
\RightLabel{\LeftR{\Weakening}}
\UnaryInf$!D, !C \fCenter !C$
\RightLabel{\RightR{\Weakening}}
\UnaryInf$!D, !C \fCenter !E, !C$
\DisplayProof
\]
Sakabehe, kita duwe !!a{proof} ing \Log{G1c} kanthi wujud iki:
\[
\Axiom$!C \fCenter !C$
\RightLabel{\LeftR{\Weakening}}
\UnaryInf$!D, !C \fCenter !C$
\RightLabel{\RightR{\Weakening}}
\UnaryInf$!D, !C \fCenter !E, !C$
\Axiom$!D \fCenter !D$
\RightLabel{\LeftR{\Weakening}}
\UnaryInf$!D, !C \fCenter !D$
\RightLabel{\RightR{\Weakening}}
\UnaryInf$!D, !C \fCenter !E, !D$
\RightLabel{\RightR{\land}}
\BinaryInf$!D, !C \fCenter !E, !C \land !D$
\Axiom$!E \fCenter !E$
\RightLabel{\LeftR{\Weakening}}
\UnaryInf$!C, !E \fCenter !E$
\RightLabel{\LeftR{\Weakening}}
\UnaryInf$!D, !C, !E \fCenter !E$
\RightLabel{\LeftR{\lif}}
\BinaryInf$ !D, !C, (!C \land !D) \lif !E \fCenter !E$
\RightLabel{\RightR{\lif}}
\UnaryInf$!C, (!C \land !D) \lif !E \fCenter !D \lif !E$
\RightLabel{\RightR{\lif}}
\UnaryInf$(!C \land !D) \lif !E \fCenter !C \lif (!D \lif !E)$
\DisplayProof
\]
Struktur bukti iki, mligine dhuwure, ora gumantung marang
!!{formula}s $!C$, $!D$, lan~$!E$ sing digunakake. Tandha-tandha
iki ing !!{proof} ndhuwur bisa diganti nganggo !!{formula}s apa wae,
lan asile tetep !!{proof} ing~\Log{G1c}. Bukti ing \Log{G1c}
iki \emph{skematis}.

Ing \Log{G3c}, aksioma identitas minangka sekuen awujud
$!A, \Gamma \Sequent \Delta, !A$ kanthi $!A$ kudu atomik.
Dadi sanajan $!D, !C \Sequent !E, !C$ katon kaya aksioma
\Log{G3c} ($\Gamma = \{!D\}$; $\Delta = \{!E\}$), formula
sing padha ing loro sisih iki njamin sekuen kasebut
\emph{pancen} aksioma yen $!C$ atomik. Pecahan bukti kita
mesthi dadi bukti jangkep ing \Log{G3c} yen $!C$, $!D$,
lan $!E$ minangka !!{formula}s atomik.

Sanajan kita durung, kanthi teges sing ketat, nuduhake yen
\[\Log{G3c} \Proves (!C \land !D) \lif !E \fCenter !C \lif (!D \lif
!E),\]
ing praktik iki kabeh sing perlu (lan bisa) kita tindakake.
Sebabe, saben sekuen awujud $!A, \Gamma \Sequent \Delta, !A$
bisa dibuktekake ing \Log{G3c}, sanajan $!A$ ora atomik.
Iki bisa dituduhake kanthi induksi tumrap jerone~$\depth{!A}$
saka~$!A$.

\begin{defn}\ollabel{defn:depth}
\emph{Jero}~$\depth{!A}$ saka !!a{formula} didhefinisikake kaya iki:
\begin{enumerate}
  \item Yen $!A$ atomik, $\depth{!A} = 0$.
  \item Yen $!A \ident \lnot !B$, $!A \ident \lforall[x][!B]$, utawa
  $!A \ident \lexists[x][!B]$, mula $\depth{!A} = \depth{!B} + 1$.
  \item Yen $!A \ident (!B \land !C)$, $(!B \lor !C)$, utawa
  $(!B \lif !C)$, mula $\depth{!A} = \max(\depth{!B},\depth{!C}) + 1$.
\end{enumerate}
\end{defn}

Ora angel mbuktekake yen $\depth{A(x)} = \depth{A(t)}$ kanggo
saben terma~$t$. Kasunyatan wigati tumrap kita yaiku ing saben
aturan logis, jerone !!{formula} utama luwih gedhe tinimbang
jerone komponen sejati, kalebu instans kuantor; formula utama
sing tetep ana ora kalebu !!{formula}s aktif sing dipecah iki.
Upamane, kita duwe:
\begin{align*}
\depth{P(x)} = \depth{P(a)} & = 0,\\
\depth{\lforall[x][P(x)]} & = 1,\\
\depth{(\lforall[x][P(x)] \lif P(a))} & = \max(1,0) + 1 = 2.
\end{align*}
Iki bisa digunakake kanggo mbuktekake asil kanthi induksi
tumrap~$\depth{!A}$, kayata asil iki.

\begin{prop}\ollabel{prop:G1c-axioms}
Kanggo saben~$!A$, ana !!a{proof} ing $\Log{G1c}$ kanggo
$!A \Sequent !A$ sing kabeh aksiomane atomik.
\end{prop}

\begin{proof}
  Kanthi induksi tumrap~$\depth{!A}$. Yen $\depth{!A} = 0$,
  $!A$ atomik lan $!A \Sequent !A$ wis dadi !!{proof}
  ing~\Log{G1c} kanthi wujud sing dibutuhake.
  Yen $\depth{!A} > 0$, $!A$ kawangun saka !!{formula}s
  sing luwih cethek, upamane $!A \ident \lnot !B$,
  $!A \ident (!B \land !C)$, utawa $!A \ident \lforall[x][!B(x)]$.
  Hipotesis induksi yaiku kanggo kabeh $!D$ kanthi
  $\depth{!D} < \depth{!A}$,
  $\Log{G1c} \Proves !D \Sequent !D$ saka aksioma atomik.

Yen $!A \ident \lnot !B$, mula $\depth{!B} < \depth{!A}$,
lan miturut hipotesis induksi ana !!a{proof}~$\pi_1$
ing \Log{G1c} saka aksioma atomik kanggo $!B \Sequent !B$.
Iki bisa kita tambahi dadi !!a{proof} kanggo
$\lnot !B \Sequent \lnot !B$ kaya iki:
\[
  \AxiomC{}
  \RightLabel{$\pi_1$}
  \Deduce$!B \fCenter !B$
  \RightLabel{\LeftR{\lnot}}
  \UnaryInf$\lnot !B, !B \fCenter $
  \RightLabel{\RightR{\lnot}}
  \UnaryInf$\lnot !B \fCenter \lnot !B$
  \DisplayProof
\]

Yen $!A \ident (!B \land !C)$, mula $\depth{!B} < \depth{!A}$
lan $\depth{!C} < \depth{!A}$. Miturut hipotesis induksi,
ana !!{proof}s~$\pi_1$ lan~$\pi_2$ ing \Log{G1c} saka aksioma
atomik kanggo $!B \Sequent !B$ lan $!C \Sequent !C$.
Iki bisa ditambahi dadi !!a{proof} kanggo
$!B \land !C \Sequent !B \land !C$ kaya iki:
\[
  \AxiomC{}
  \RightLabel{$\pi_1$}
  \Deduce$!B, \fCenter !B$
  \RightLabel{\LeftR{\Weakening}}
  \UnaryInf$!B, !C, \fCenter !B$
  \RightLabel{\LeftR{\land}}
  \UnaryInf$!B \land !C, !C\fCenter !B$
  \RightLabel{\LeftR{\land}}
  \UnaryInf$!B \land !C, !B \land !C \fCenter !B$
  \AxiomC{}
  \RightLabel{$\pi_2$}
  \Deduce$!C, \fCenter !C$
  \RightLabel{\LeftR{\Weakening}}
  \UnaryInf$!B, !C, \fCenter !C$
  \RightLabel{\LeftR{\land}}
  \UnaryInf$!B \land !C, !C \fCenter !C$
  \RightLabel{\LeftR{\land}}
  \UnaryInf$!B \land !C, !B \land !C \fCenter !C$
  \RightLabel{\RightR{\land}}
  \BinaryInf$!B \land !C, !B \land !C \fCenter !B \land !C$
  \RightLabel{\LeftR{\Contraction}}
  \UnaryInf$!B \land !C \fCenter !B \land !C$
  \DisplayProof
  \]

Saiki upamane $!A \ident \lforall[x][!B(x)]$.
Pilih !!a{constant}~$c$ sing ora ana ing $\lforall[x][!B(x)]$.
Mula $\depth{!B(c)} = \depth{!B(x)} < \depth{!A}$,
lan miturut hipotesis induksi ana !!a{proof}~$\pi_1$
ing \Log{G1c} saka aksioma atomik kanggo
$!B(c) \Sequent !B(c)$. Iki bisa ditambahi dadi
!!a{proof} kanggo $\lforall[x][!B(x)] \Sequent \lforall[x][!B(x)]$
kaya iki:
\[
  \AxiomC{}
  \RightLabel{$\pi_1$}
  \Deduce$!B(c) \fCenter !B(c)$
  \RightLabel{\LeftR{\lforall}}
  \UnaryInf$\lforall[x][!B(x)] \fCenter !B(c)$
  \RightLabel{\RightR{\lforall}}
  \UnaryInf$\lforall[x][!B(x)] \fCenter \lforall[x][!B(x)]$
  \DisplayProof
\]

Kasus sing isih kari ditinggal minangka latihan.
\end{proof}

\begin{prob}
  Jangkepana bukti \olref{prop:G1c-axioms} kanthi nindakake
  kasus nalika $!A \ident (!B \lor !C)$, $!A \ident (!B \lif !C)$,
  lan $!A \ident \lexists[x][!B(x)]$.
\end{prob}

Kanggo ancas kita, kita uga bisa nggunakake !!{proof}
\[
\AxiomC{}
\RightLabel{$\pi_1$}
\Deduce$!B \fCenter !B$
\RightLabel{\RightR{\lnot}}
\UnaryInf$\fCenter !B, \lnot !B$
\RightLabel{\LeftR{\lnot}}
\UnaryInf$\lnot !B \fCenter \lnot !B$
\DisplayProof
\]
ing kasus nalika $!A \ident \lnot !B$. Nanging ana sistem
sekuen sing ora ngidini cara iki. Kalkulus sekuen intuisionistik
\Log{G1i} padha karo \Log{G1c}, nanging saben sekuen diwatesi
supaya ngemot paling akeh siji !!{formula} ing sisih tengen.
Yen $\Delta = \emptyset$, !!{proof} kapisan uga dadi
!!a{proof} ing~\Log{G1i}, nanging kapindho ora, amarga
ana sekuen sing ngemot $!B$ lan~$\lnot !B$ ing tengen.

Gatekna yen sanajan ing \Log{G1c}, kita ora bisa
nggunakake urutan aturan walikan ing kasus nalika
$!A \ident \lforall[x][!B(x)]$. Iki \emph{dudu} !!a{proof}
\[
  \AxiomC{}
  \RightLabel{$\pi_1$}
  \Deduce$!B(c) \fCenter !B(c)$
  \RightLabel{\RightR{\lforall}}
  \UnaryInf$!B(c) \fCenter \lforall[x][!B(x)]$
  \RightLabel{\LeftR{\lforall}}
  \UnaryInf$\lforall[x][!B(x)] \fCenter \lforall[x][!B(x)]$
  \DisplayProof
\]
amarga syarat eigenvariabel kanggo~$\RightR{\lforall}$ dilanggar.

\begin{prop}\ollabel{prop:G3c-axioms}
Saben sekuen awujud $!A, \Gamma \Sequent \Delta, !A$
(kanthi $!A$ \emph{ora} kudu atomik) bisa dibuktekake
ing~\Log{G3c}.
\end{prop}

\begin{proof}
Buktine padha banget karo bukti \olref{prop:G1c-axioms},
nanging kita kudu ngati-ati tumrap hipotesis induksi.
Kasus nalika $n = \depth{!A} = 0$ gampang maneh.
Yen $n=0$, $!A$ atomik, lan
$!A, \Gamma \Sequent \Delta, !A$ minangka aksioma~\Log{G3c}.

Saiki upamane $n>0$. Kita misahake kasus maneh.
Hipotesis induksi yaiku kanggo kabeh $!D$ kanthi
$\depth{!D} < n$, lan kabeh $\Pi$ lan $\Lambda$,
$\Log{G3c} \Proves !D, \Pi \Sequent \Lambda, !D$.
Iki kasus kanggo $!A \ident (!B \land !C)$.
Kita nggunakake hipotesis induksi kaping pindho:
sapisan kanggo $!D = !B$, $\Pi = !C, \Gamma$, lan
$\Lambda = \Delta$, kanggo oleh !!a{proof}~$\pi_1$
kanggo $!B, !C, \Gamma \Sequent \Delta, !B$;
lan sapisan kanggo $!D = !C$, $\Pi = !B, \Gamma$,
lan $\Lambda = \Delta$, kanggo oleh !!a{proof}~$\pi_2$
kanggo $!B, !C, \Gamma \Sequent \Delta, !C$.
Iki !!{proof} kanggo~$!B \land !C, \Gamma \Sequent \Delta, !B \land !C$:
\[
\AxiomC{}
\RightLabel{$\pi_1$}
\Deduce$!B, !C, \Gamma \fCenter \Delta, !B$
\RightLabel{\LeftR{\land}}
\UnaryInf$!B \land !C, \Gamma \fCenter \Delta, !B$
\AxiomC{}
\RightLabel{$\pi_2$}
\Deduce$!B, !C, \Gamma \fCenter \Delta, !C$
\RightLabel{\LeftR{\land}}
\UnaryInf$!B \land !C, \Gamma \fCenter \Delta, !C$
\RightLabel{\RightR{\land}}
\BinaryInf$!B \land !C, \Gamma \fCenter \Delta, !B \land !C$
\DisplayProof
\]

Kanggo $!A \ident \lforall[x][!B(x)]$, pilih
!!a{constant} sing ora ana ing~$\lforall[x][!B(x)]$,
$\Gamma$, utawa~$\Delta$. Trapna hipotesis induksi
kanggo $!D = !B(c)$, $\Pi = \lforall[x][!B(x)], \Gamma$,
lan $\Lambda = \Delta$. Iki menehi !!a{proof}~$\pi_1$
kanggo $!B(c), \lforall[x][!B(x)], \Gamma \Sequent \Delta, !B(c)$,
lan saka bukti iki kita oleh:
\[
\AxiomC{}
\RightLabel{$\pi_1$}
\Deduce$!B(c), \lforall[x][!B(x)], \Gamma \fCenter \Delta, !B(c)$
\RightLabel{\LeftR{\lforall}}
\UnaryInf$\lforall[x][!B(x)], \Gamma \fCenter \Delta, !B(c)$
\RightLabel{\RightR{\lforall}}
\UnaryInf$\lforall[x][!B(x)], \Gamma \fCenter \Delta, \lforall[x][!B(x)]$
\DisplayProof
\]
Syarat eigenvariabel kanggo $\RightR{\lforall}$
wis dicukupi kanthi pilihan~$c$ mau.
\end{proof}

\begin{prob}
  Terusna bukti \olref{prop:G3c-axioms} kanthi nindakake
  kasus nalika $!A \ident (!B \lif !C)$ lan
  $!A \ident \lexists[x][!B(x)]$.
\end{prob}

\end{document}
